\documentclass[10pt,twocolumn,letterpaper]{article}

% PROVISIONAL CVPR REVIEW FORMAT.
% The active cvpr.sty is byte-identical to the pinned official
% cvpr-org/author-kit snapshot recorded in CVPR_TEMPLATE_PROVENANCE.md.
% That upstream snapshot identifies itself as CVPR 2026 and is used
% provisionally until an official CVPR 2027 author kit is published.
\usepackage[review]{cvpr}

\usepackage{multirow}
\usepackage{microtype}

\definecolor{cvprblue}{rgb}{0.21,0.49,0.74}
\usepackage[pagebackref,breaklinks,colorlinks,allcolors=cvprblue]{hyperref}

\def\paperID{*****}
\def\confName{CVPR}
\def\confYear{2027}

\title{When Should SAM 3 Write to Memory?\\
A Controlled Study of Signal-Informed Memory Admission for
Video Object Segmentation}

\author{Anonymous CVPR 2027 Submission}

\renewcommand{\thesection}{S\arabic{section}}
\renewcommand{\thesubsection}{S\arabic{section}.\arabic{subsection}}
\renewcommand{\thetable}{S\arabic{table}}

\begin{document}
\maketitlesupplementary

This supplement expands the frozen implementation, evaluation, and
reproducibility details used by the main paper.  It does not introduce new
confirmatory analyses or modify the one-touch TEST conclusions.

\section{Frozen substrate and evaluation boundary}
\label{sec:supp-substrate}

The final substrate is frozen SAM 3 VOS/PVS constructed with
\texttt{build\_sam3\_video\_model()} at SAM source commit
\texttt{8f0b7f4d4e7eda2ed606ebde6702c93359ad01da}.
No SAM gradient, LoRA, fine-tuning, detector modification, or new SAM
architecture is used.  The gate parameter budget is at most 50K learned
parameters.

All final experiments use an NVIDIA RTX 4080 SUPER with 16 GB VRAM under
WSL2 Ubuntu 22.04.5, Python 3.12.13, PyTorch 2.10.0+cu128, and CUDA 12.8.
The realized Python environment is recorded by
\texttt{requirements.lock.txt}, whose frozen SHA256 is
\texttt{c943caf335003594\allowbreak{}dc93d9042267695e\allowbreak{}70c0e955652257ca\allowbreak{}ae8a3ff11a0efa6a}.

MOSEv2 train contains 3,666 videos in the available release used here.
The released validation split is not used for the endpoint construction
because it does not provide the required dense per-frame annotation.
Video is the statistical cluster.  HARD\_TEST80 contains 80 videos and
506 qualifying primary events; REPRESENTATIVE\_TEST40 contains 40 videos
and 76 qualifying primary events.

TEST prediction-derived evidence is touched exactly once after the final
freeze.  No TEST threshold search or threshold retuning is permitted, and
unsupported development write-rate targets are not interpolated post hoc on
TEST.

\section{Gate architecture and routing}
\label{sec:supp-gates}

B0 is the native ungated write reference.  B1 is a fixed manual
quality--temporal rule.  B2 is the learned quality--temporal condition.
B3-S adds native self-identity information, B3-R adds competitor-relative
identity information, and B5 is the DMS-lite write-side comparator used in
the main study.

\begin{table*}[t]
\centering
\small
\caption{Frozen signal ladder and headline operating thresholds.  Parameter
counts refer only to learned gate parameters; SAM 3 remains frozen.}
\label{tab:supp-gates}
\begin{tabular}{llrr}
\toprule
Policy & Information available to admission & Learned params & $\tau$ \\
\midrule
B0   & Native ungated reference & 0 & -- \\
B1   & Manual quality + temporal & 0 & 0.10 \\
B2   & Five quality/temporal features & 4,994 & 0.10 \\
B3-S & B2 + native self-pointer similarity & 5,122 & 0.20 \\
B3-R & B3-S + competitor-relative pointer similarity & 5,250 & 0.20 \\
B5   & Confidence/presence reliability comparator & 0 & 0.70 \\
\bottomrule
\end{tabular}
\end{table*}

B2 uses two independent failure-typed MLP heads, one for drift and one for
theft.  Each head has the architecture
$5\!\rightarrow\!64\!\rightarrow\!32\!\rightarrow\!1$ with ReLU hidden
activations, for 2,497 parameters per head and 4,994 learned parameters in
total.  Training is deterministic full-batch PyTorch CPU float64 with
independently weighted binary cross-entropy losses, AdamW learning rate
0.001, weight decay 0.0001, betas $(0.9,0.999)$, epsilon $10^{-8}$, and
1,000 epochs without early stopping or DEV tuning.  The TRAIN-only
normalizer is a z-score transform fitted on finite B2 training rows.

B3-S changes the input dimension to six by adding
\texttt{ptr\_sim\_anchor\_fp32}; each failure head is
$6\!\rightarrow\!64\!\rightarrow\!32\!\rightarrow\!1$, for 5,122 learned
parameters in total.  B3-R adds
\texttt{max\_comp\_anchor\_cos\_fp32}; each failure head is
$7\!\rightarrow\!64\!\rightarrow\!32\!\rightarrow\!1$, for 5,250 learned
parameters in total.

Identity cosine calculations use FP32 with autocast disabled, CUDA TF32
disabled, and \texttt{torch.set\_float32\_matmul\_precision("highest")}.
The availability mask is
\texttt{pointer\_valid = object\_score\_logit > 0}; it is used only for
routing and is not supplied as a predictive numeric feature.  B3-S falls
back to B2 when self identity is unavailable.  B3-R routes to B3-R when
relational identity is available, to B3-S when only self identity is
available, and to B2 when self identity is unavailable.

For B2, B3-S, and B3-R, any non-finite required B2 base feature fails closed
with object admission score zero.  Independent drift and theft unsafe
probabilities are converted to safe probabilities and combined
conservatively by their minimum.  Object scores are then aggregated by the
minimum across tracked objects before the thresholded frame-level
\textsc{Admit}/\textsc{Block} decision.

\section{Final HARD\_TEST80 operating-point results}
\label{sec:supp-hard}

\begin{table*}[t]
\centering
\small
\caption{Frozen HARD\_TEST80 headline results.  Peak VRAM and runtime are
reported for the executed condition and are descriptive engineering
measurements rather than statistical endpoints.}
\label{tab:supp-hard}
\begin{tabular}{lrrrrr}
\toprule
Policy & POR@30 & ITR@30 & Write rate & Peak VRAM (GB) & Runtime (s) \\
\midrule
B0   & 0.7411 & 0.0593 & 1.0000 & 15.81 & 2572.7 \\
B1   & 0.7273 & 0.0514 & 0.3199 & 13.11 & 1852.5 \\
B2   & 0.7569 & 0.0316 & 0.3233 & 13.06 & 1878.9 \\
B3-S & 0.7648 & 0.0395 & 0.2906 & 13.06 & 1878.1 \\
B3-R & 0.7648 & 0.0336 & 0.2923 & 13.06 & 1888.7 \\
B5   & 0.7530 & 0.0336 & 0.3080 & 13.03 & 1845.7 \\
\bottomrule
\end{tabular}
\end{table*}

The preregistered primary contrast is B3-S minus B2 on HARD\_TEST80.  The
observed POR@30 difference is $0.007905$, with video-clustered BCa 95\%
confidence interval $[-0.002037,\,0.020882]$.  The realized write rates are
0.290565 for B3-S and 0.323269 for B2, an absolute difference of 0.032703.
Because this exceeds the frozen 0.02 tolerance, the primary contrast is
classified as \textsc{Rate-Mismatch} and does not receive a matched-rate
primary interpretation.

The sole primary endpoint is POR@30.  ITR@30 is secondary.  Paired
video-clustered BCa 95\% intervals use 50,000 bootstrap replicates, seed 52,
midrank bias correction, and delete-one-video jackknife acceleration.

\section{Representative TEST cohort}
\label{sec:supp-representative}

\begin{table}[t]
\centering
\small
\caption{REPRESENTATIVE\_TEST40 descriptive results.}
\label{tab:supp-representative}
\begin{tabular}{lrrr}
\toprule
Policy & POR@30 & ITR@30 & Write rate \\
\midrule
B0   & 0.5789 & 0.0000 & 1.0000 \\
B1   & 0.6711 & 0.0132 & 0.3375 \\
B2   & 0.7237 & 0.0132 & 0.3739 \\
B3-S & 0.7237 & 0.0132 & 0.3581 \\
B3-R & 0.7237 & 0.0000 & 0.3652 \\
B5   & 0.6974 & 0.0132 & 0.3977 \\
\bottomrule
\end{tabular}
\end{table}

These 40 videos and 76 qualifying events provide the pre-specified
external-validity cohort.  They are reported descriptively and do not replace
the HARD\_TEST80 confirmatory inference.

\section{Exact-budget neutral controls and supported-rate discipline}
\label{sec:supp-rate-controls}

For every signal-based condition, the neutral counterpart receives exactly
the same number of admitted writes in each video.  The admitted frame
identities are selected independently of the signal gate while the count is
held fixed.  Signal and neutral admit counts are asserted to match before the
paired diagnostic is accepted.  These controls therefore separate the write
budget from the signal-based timing choice.

Write-rate operating points are selected from frozen DEV support before
TEST.  Unsupported development targets are not interpolated or repaired
after TEST, and the evaluator records
\texttt{threshold\_search\_performed = False} and
\texttt{threshold\_retuning\_performed = False}.

\section{Post-defense full-video segmentation probe}
\label{sec:supp-exp055}

EXP055 is supplementary, post-defense, and nonconfirmatory.  It evaluates
ordinary full-video segmentation quality for frozen B0 versus frozen B2 on
80 outcome-independently sampled untouched videos.  No retraining, threshold
search, or threshold retuning is performed.

\begin{table*}[t]
\centering
\small
\caption{EXP055 full-video mIoU results.  Confidence intervals are
video-clustered BCa 95\% intervals for B2 minus B0.}
\label{tab:supp-exp055}
\begin{tabular}{lrrrr}
\toprule
Estimator & B0 & B2 & B2$-$B0 & 95\% CI \\
\midrule
Pooled object-frame mean IoU
  & 0.736943 & 0.678104 & $-0.058839$ & $[-0.118452,\,0.001673]$ \\
Video-balanced mean IoU
  & 0.704008 & 0.666728 & $-0.037280$ & $[-0.099020,\,0.020909]$ \\
\bottomrule
\end{tabular}
\end{table*}

The analysis contains 9,778 evaluable GT-visible non-conditioning
object-frames and uses 50,000 bootstrap replicates.  B2 uses its frozen
threshold $\tau=0.1$ and realizes pooled write rate 0.466350.  Peak VRAM is
11.6997 GB and total tracker runtime is 2149.646 s.  Both confidence
intervals include zero, so this supplementary probe establishes neither a
statistically reliable full-video mIoU improvement nor degradation for B2
and does not modify the frozen confirmatory TEST conclusions.

\section{Reproducibility artifacts}
\label{sec:supp-repro}

The public-release environment record is
\texttt{docs/ENVIRONMENT.md}; \texttt{environment.yml} records the Conda
environment specification; \texttt{requirements.lock.txt} records the
realized Python package environment; and \texttt{docs/SAM3\_INSTALL.md}
records the SAM source pin and installation-specific dependency fixes.

The frozen confirmatory protocol is recorded in
\texttt{PREREGISTRATION.md}.  Experiment commands and provenance are tracked
in \texttt{EXPERIMENT\_REGISTRY.md}, while the chronological current-state
record is \texttt{RESEARCH\_STATE.md}.  The EXP053 final TEST directory
contains \texttt{artifact\_manifest.json}, which hashes the final config,
script, and major output files.

The frozen learned-model artifacts are:
\begin{itemize}
\item B2:
  \path{experiments/EXP029_b2core_train/model.json}, SHA256
  \texttt{6160a6be9da20588\allowbreak{}08c16182abe03443\allowbreak{}014806273df46fe5\allowbreak{}9a86411ffc869ecf};
\item B3-S/B3-R:
  \path{experiments/EXP031_b3_train/model.json}, SHA256
  \texttt{235076b86aa0975b\allowbreak{}3ec624aae703575f\allowbreak{}d4f030381b6af400\allowbreak{}8c3808f2db71847a}.
\end{itemize}

Historical planning artifacts remain available for provenance, but later
frozen amendments and executed experiment records supersede them where the
repository explicitly records a methodological correction.

\end{document}
